% Reconstructed text from V5 PDF page 13; non-native source
\par\smallskip\noindent\rule{\textwidth}{0.4pt}\par{\sffamily\footnotesize Original V5 page 13 -- recovered text}\par\begin{flushleft}\footnotesize\ttfamily
\par\smallskip
\par\smallskip
\noindent Diagnostics and Serving Infrastructure for MAS 13\par
\par\smallskip
\par\smallskip
\par\smallskip
\noindent  Role), and three edge types connecting them (§4.1). The system backtracks through this\par
\noindent  graph in reverse topological order using synthesized virtual oracles Ok = ⟨Goal, Pre-\par
\noindent  conditions, Key Evidence, Acceptance Criteria⟩, then applies counterfactual attribution\par
\noindent  through four progressive stages: local analysis, planning-control analysis, data-flow anal-\par
\noindent  ysis, and deviation-aware analysis (§4.2–4.3). A RAG knowledge base constructed from\par
\noindent  165 GAIA and 33 AssistantBench examples supports subtask decomposition.\par
\noindent  Using DeepSeek-V3.2 (671B parameters, approximately 37B active via sparse MoE) as\par
\noindent  its backbone, CHIEF achieves the highest reported agent-level accuracy at 77.59\% on\par
\noindent  handcrafted scenarios and 76.80\% on algorithm-generated scenarios (Table 1). Step-level\par
\noindent  accuracy is 29.31\% (handcrafted) and 52.00\% (algorithm-generated). The system con-\par
\noindent  sumes 55,085 tokens per trajectory on handcrafted scenarios and 19,504 on algorithm-\par
\noindent  generated ones (Table 2). Reported costs are presented in original units; we avoid cross-\par
\noindent  paper cost ratios when workloads and measurement conditions differ. The method was\par
\noindent  tested on TITAN RTX with 32GB RAM using API inference (§5). Its effectiveness de-\par
\noindent  pends on graph fidelity: hallucinated edges in the causal graph propagate attribution errors\par
\noindent  through the backtracking process (§7). A single-benchmark evaluation (Who\&When) and\par
\noindent  single-root-cause assumption (no cumulative error propagation verification) limit gener-\par
\noindent  alizability.\par
\noindent  • TrajAudit [35] addresses the challenge of long execution trajectories (often thousands of\par
\noindent  steps in coding agents) through two preprocessing stages before LLM-based diagnosis.\par
\noindent  Semantic Structure Folding (SSF) merges consecutive actions with similar semantics,\par
\noindent  reducing the trajectory length while preserving the decision structure; the paper reports\par
\noindent  94.6\% of steps folded on average (§VII-D). Progressive Failure Reasoning (PFR) then\par
\noindent  applies a multi-round narrowing strategy that progressively identifies the earliest decisive\par
\noindent  error step.\par
\noindent  The system was tested with five LLM backbones: Claude Sonnet 4.5 (default), Gemini\par
\noindent  3.1 Pro, GPT-5.2, DeepSeek V4 Pro, and Qwen3-235B (Table VI). The best configuration\par
\noindent  achieves 52.0\% step-level RootSE accuracy without reference trajectories and 49.0\% with\par
\noindent  reference trajectories, representing +21.6 and +10.8 percentage-point gains over direct-\par
\noindent  prompting baselines respectively (Table II). Token consumption averages 138K per case\par
\noindent  (Table VII, §VII-C), making TrajAudit the most token-intensive single-trajectory method\par
\noindent  in the survey.\par
\noindent  Several limitations are notable (§IX). The dataset was filtered from 500 to 102 instances,\par
\noindent  excluding ambiguous and multi-failure cases. The method is evaluated on single-agent\par
\noindent  coding tasks only; 37\% of MAS failures are inter-agent cooperation breakdowns, which\par
\noindent  TrajAudit does not address. A strong positional bias exists: accuracy drops to 19–26\%\par
\noindent  when errors occur in the first 20\% of steps, and 81\% of mispredictions are later than the\par
\noindent  true error step by an average of 14.5 steps. Ablation confirms SSF as critical (−11.4\%\par
\noindent  vs. −3.5\% for PFR, Table VII).\par
\noindent  • AgentRx [3] synthesizes behavioral constraints from toolset schemas and domain poli-\par
\noindent  cies, then uses an LLM judge to evaluate each step against the generated constraints.\par
\noindent  The method operates in two phases: first, constraint generation produces semantic asser-\par
\noindent  tions for each step based on the available tool definitions and optional domain policy Π\par
\noindent  (9 failure categories in a structured taxonomy); second, a judge protocol evaluates each\par
\noindent  step against these constraints. On τ -bench [42] (airline and retail scenarios), AgentRx\par
\noindent  achieves 42.7\% step accuracy, with the AgentRx W\&W* variant showing +23.9 percent-\par
\noindent  age points over the 18.8\% baseline (Table 3). Cost is \$0.18 per trajectory versus \$0.05 for\par
\noindent  judge-only evaluation (3.6× overhead), with constraint generation consuming 81\% of the\par
\end{flushleft}
